\section{Experiment Setup}
\todo{Fill in numbers of generated data}
To test our benchmark we generate a SolidBench dataset with default scale = $0.1$, resulting in x triples over y pods.
For the query sequence generation, we generate 6 sequences with in total $136$ query instantiations. The hyperparameters used for the generation are outlined in Table \ref{tab:hyperparameters}.

The purpose of our experiment is two-fold, one is to evaluate the characteristics of the simulated query sequences. 
Specifically, the impact our three benchmark design choices have on the correlations between queries. 
For caching-based algorithms, this correlation is primarily measured in the \textit{hit rate} and \textit{evictions}, thus to evaluate our benchmark we will investigate these metrics.
The other is to evaluate the actual effectiveness of the implemented caching approaches in a single-threaded link traversal engine.
While caching is not a new concept, the introduced caching approaches have not been rigorously tested in the single threaded execution environment of Comunica.
Thus, we can immediately use our proposed benchmark to investigate the performance characteristics of our proposed caching.

\subsection{Cache Hit Rate Experiment}
To investigate the theoretical performance benefits of caching in link traversal, we additionally run experiments with synthetically set cache hit rates.

\todo{Explain the experiment}

This experiment will not only provide new insights into the possible benefits of caching approaches as different levels of performance, it will also show that for a given average hitrate in a sequence the overhead of evictions and cache management costs, as this cost is excluded from the synthetic hitrate experiments.
Thus, giving a more complete view of the performance landscape of such approaches and the realism of our proposed benchmark.
